% AUTHOR: Amlal El Mahrouss
% PURPOSE: DN001.10: Notes for Several Formulas.

\documentclass[11pt, a4paper]{article}
\usepackage{graphicx}
\usepackage{listings}
\usepackage{amsmath,amssymb,amsthm}
\usepackage{xcolor}
\usepackage{hyperref}
\usepackage{mathtools}
\usepackage[margin=0.5in,top=1in,bottom=1in]{geometry}

\title{Third Notes for Several Formulas.}
\author{Amlal El Mahrouss\\\texttt{amlal@nekernel.org}}
\date{February 21, 2026}

\begin{document}

\maketitle

\begin{center}
	\rule[0.01cm]{17cm}{0.01cm}
\end{center}

\begin{abstract}
\begin{center}
    We'll develop new formulas in this paper, where we will cover some new formulas and properties.
\end{center}
\end{abstract}

\section{DEFINITIONS}

We assume the following to be true, and $U(t), N(u)$ to be assumed defined:

\section{DEFINITION OF A G-FUNCTION:}

We define $G(\alpha, \Phi, n)$ as:

\begin{equation}
    G(\alpha, \Phi, n) \coloneqq \frac{\partial U(\alpha)}{\partial \alpha} = \sum_{k=N(\Phi)}^{n} U(|\alpha|)
\end{equation}
For $\Phi \geq 1, \forall \alpha, \Phi \in \mathbb{R}, \forall n, N(\Phi) \in \mathbb{Z}$

\section{DEFINITION OF A U-FUNCTION:}

Let an U-Function in the form of:

\begin{equation}
    U(x) \in \mathbb{Z}
\end{equation}
Be $U(x) > 0$.

\section{DEFINITION OF A N-FUNCTION:}

Let an N-Function in the form of:

\begin{equation}
    N(x) \in \mathbb{R}
\end{equation}
Be $N(x) > 0$.

\end{document}
